


\section{A complete rate, covenant and ownership path}
An exact combined example shows why the financing steps must be dated. Two
possible companies each start with cash 16 and loan face 100. One pays a
fixed coupon of 5; the other pays 7 after the specified floating-rate reset.
Each then makes an operating payment of 6. A covenant tests cash of at least
5 immediately after these payments and before financing. Breach forbids new
borrowing but permits an equity cure. Management requires closing cash 8.
These are assumed contract terms, not a claim about a standard covenant.

The fixed-rate company has cash $16-5-6=5$, meets the covenant and needs 3.
A candidate debt issue with maximum face 5, conditional solvency probability
$4/5$, recovery 2 and gross risk-free return $21/20$ has proceeds $88/21$.
Holding that quoted per-face price fixed for this conditional stress, issuing
face $315/88$ raises exactly 3. The floating-rate company has cash
$16-7-6=3$, breaches the covenant and needs 5. Debt is unavailable; even an
improper borrowing override at the maximum face would raise only $88/21<5$.

Suppose the permitted equity cure has fee 1, pre-issue economic equity value
20 and 100 incumbent shares. A gross subscription of 6 provides net cash 5.
The post-issue equity value is 25, new investors own $6/25$, and the common
share price is $19/100$. Their $600/19$ new shares have value 6, while
incumbents retain value 19. Both companies finish with cash 8, but their new
debt and ownership differ. Merging them on closing cash would conceal those
differences. Testing the covenant only after the cure would erase the earlier
breach and wrongly admit the borrowing alternative.

This example holds quoted financing prices and the pre-issue equity valuation
fixed to expose event ordering and allocation. It does not solve the feedback
from financing to default and prices. A source correction introduces another
interaction: withdrawing an unsupported fixed-rate classification can reopen
the floating-rate history, invalidating an earlier single-history result.
If management later receives a rate notice, its admissible policies can change
even when the analyst's reported balances do not. Shared computations must
therefore depend on evidence version, actor information, contracts and the
requested outcomes, as well as current cash.

\section{Executing promises, carrying amounts and ownership}
The combined example can be executed through the existing typed company
runtime once its state distinguishes contractual face from carrying amount.
For the fixed-rate history, the authored initial-measurement convention posts
cash of 3 and a liability carrying amount of 3, while contractual face increases
by $315/88$. The closing liability carrying amount is 103 and outstanding face
is $9115/88$. Subsequent effective-interest measurement and settlement require
further methods; substituting carrying amount for the promise would lose the
very exposure the example is intended to preserve.

For the floating-rate history, subscription proceeds of 6 increase cash and
contributed capital; the direct issue cost of 1 reduces both. The typed state
retains 100 incumbent shares and $600/19$ new shares. Under this explicitly
synthetic posting convention, both histories report assets of 112. The fixed
history has liabilities of 103 and book equity of 9; the floating history has
liabilities of 100 and book equity of 12. These book amounts differ from the
economic equity values used to price the subscription. No reporting-standard
admission follows from the balanced journals.

The reference adapter executes coupon payment, operations, covenant observation
and financing in that order, using the unchanged common company executor.
It rejects borrowing after the floating-rate breach, proceeds inconsistent
with the supplied per-face quote, a zero issue price, duplicate issuance and
payment without sufficient cash. A mutation that omits the share-count update
still produces balanced accounts and completes the accounting path; the
independent ownership comparison detects the error. Accounting conservation
alone therefore cannot validate a financing method.

The same runtime also executes the earlier information example. Starting with
cash 6, a fixed coupon 5 leaves 1, while a floating coupon 7 cannot be paid.
A hedge costing 1 caps the coupon at 5. An informed manager hedges only the
floating history and retains cash 1 or 0; a manager who must act identically in
both histories can hedge both and retain 0. Withdrawing the earlier fixed-only
evidence reopens the payment-refusal branch. The negative cash in the algebraic
shortfall calculation denotes unmet funding; the runtime refuses the payment
atomically and keeps the opening cash, rather than recording an impossible
settlement. A missed-payment event and any ensuing default process would need
separate declared methods.
